\chapter{2011年广东卷（文）}

\section{单选题}

\begin{problem}
设复数 \(z\) 满足 \(\i z = 1\)，其中 \(\i\) 为虚数单位，则 \(z =\)（\quad）

\choices
    {\(-\i\)}
    {\(\i\)}
    {\(-1\)}
    {\(1\)}
\end{problem}

\begin{answer}
A
\end{answer}

\begin{solution}
由 \(\i z=1\)，两边同乘 \(-\i\)，利用 \((-\i)\i=1\)，得 \(z=-\i\)，所以选 A.
\end{solution}

\begin{problem}
已知集合 \(A = \{(x, y) | x, y\) 为实数，且 \(x^2 + y^2 = 1\}\)，\(B = \{(x, y) | x, y\) 为实数，且 \(x + y = 1\}\)，则 \(A \cap B\) 的元素个数为（\quad）

\choices
    {\(4\)}
    {\(3\)}
    {\(2\)}
    {\(1\)}
\end{problem}

\begin{answer}
C
\end{answer}

\begin{solution}
点 \((x,y)\) 属于 \(A\cap B\) 当且仅当同时满足 \(x^2+y^2=1\) 和 \(x+y=1\). 由 \(y=1-x\)，代入得 \(x^2+(1-x)^2=1\)，即 \(2x(x-1)=0\)，所以 \(x=0\) 或 \(x=1\). 相应的点为 \((0,1)\) 和 \((1,0)\)，共有 \(2\) 个元素，故选 C.
\end{solution}

\begin{problem}
已知向量 \(\bs{a} = (1, 2)\)，\(\bs{b} = (1, 0)\)，\(\bs{c} = (3, 4)\). 若 \(\lambda\) 为实数，\((\bs{a} + \lambda \bs{b}) \parallel \bs{c}\)，则 \(\lambda = \)（\quad）

\choices
    {\( \frac{1}{4} \)}
    {\( \frac{1}{2} \)}
    {\(1\)}
    {\(2\)}
\end{problem}

\begin{answer}
B
\end{answer}

\begin{solution}
有 \(\bs{a}+\lambda\bs{b}=(1+\lambda,2)\). 两个向量平行，当且仅当其坐标行列式为零，因此
\[
4(1+\lambda)-2\cdot3=0
\]
解得 \(\lambda=\frac{1}{2}\)，故选 B.
\end{solution}

\begin{problem}
函数 \( f(x) = \frac{1}{1 - x} + \lg (1 + x) \) 的定义域是（\quad）

\choices
    {\((-\infty, -1)\)}
    {\((1, +\infty)\)}
    {\((-1,1)\cup (1, + \infty)\)}
    {\((-\infty, +\infty)\)}
\end{problem}

\begin{answer}
C
\end{answer}

\begin{solution}
分式有意义要求 \(1-x\ne0\)，即 \(x\ne1\)；对数有意义要求 \(1+x>0\)，即 \(x>-1\). 合并得定义域 \((-1,1)\cup(1,+\infty)\)，故选 C.
\end{solution}

\begin{problem}
不等式 \(2x^{2} - x - 1 > 0\) 的解集是（\quad）

\choices
    {\( \left(-\frac{1}{2},1\right) \)}
    {\((1, + \infty)\)}
    {\((-\infty, 1) \cup (2, +\infty)\)}
    {\(\left(-\infty, -\frac{1}{2}\right) \cup (1, +\infty)\)}
\end{problem}

\begin{answer}
D
\end{answer}

\begin{solution}
因式分解得 \(2x^2-x-1=(2x+1)(x-1)\)，两个根为 \(-\frac{1}{2}\) 和 \(1\). 二次项系数为正，所以积在两根外侧为正，且不取使积为零的两个端点. 因此解集为 \(\left(-\infty,-\frac{1}{2}\right)\cup(1,+\infty)\)，选 D.
\end{solution}

\begin{problem}
已知平面直角坐标系 \(xOy\) 上的区域 \(D\) 由不等式组 \(\begin{cases}
0 \leqslant x \leqslant \sqrt{2},\\
y \leqslant 2,\\
x \leqslant \sqrt{2}y
\end{cases}\) 给定. 若 \(M(x, y)\) 为 \(D\) 上的动点，点 \(A\) 的坐标为 \((\sqrt{2}, 1)\)，则 \(z = \overrightarrow{OM} \cdot \overrightarrow{OA}\) 的最大值为（\quad）

\choices
    {\(3\)}
    {\(4\)}
    {\( 3\sqrt{2} \)}
    {\(4\sqrt{2}\)}
\end{problem}

\begin{answer}
B
\end{answer}

\begin{solution}
由坐标得 \(\overrightarrow{OM}=(x,y)\)，\(\overrightarrow{OA}=(\sqrt{2},1)\)，故
\[
z=\overrightarrow{OM}\cdot\overrightarrow{OA}=\sqrt{2}x+y
\]
由 \(0\leqslant x\leqslant\sqrt{2}\) 和 \(y\leqslant2\)，有 \(z\leqslant\sqrt{2}\cdot\sqrt{2}+2=4\). 点 \((\sqrt{2},2)\) 还满足 \(\sqrt{2}\leqslant\sqrt{2}\cdot2\)，所以确实属于 \(D\)，并在该点取到上界. 最大值为 \(4\)，故选 B.
\end{solution}

\begin{problem}
正五棱柱中，不同在任何侧面且不同在任何底面的两顶点的连线称为它的对角线，那么一个正五棱柱对角线的条数共有（\quad）

\choices
    {\(20\)}
    {\(15\)}
    {\(12\)}
    {\(10\)}
\end{problem}

\begin{answer}
D
\end{answer}

\begin{solution}
同一底面的两个顶点不符合条件，所以所求线段必须分别取自上、下两个底面. 上下底面各有 \(5\) 个顶点，跨两个底面的连线共有 \(5\times5=25\) 条.

其中，连接对应顶点的 \(5\) 条线段是侧棱；每个侧面还有 \(2\) 条面内对角线，五个侧面共有 \(5\times2=10\) 条. 这 \(15\) 条线段均位于某个侧面内，不能算作所求对角线. 因此对角线条数为 \(25-5-10=10\)，故选 D.
\end{solution}

\begin{problem}
设圆 \(C\) 与圆 \(x^{2}+(y-3)^{2}=1\) 外切，且与直线 \(y=0\) 相切，则 \(C\) 的圆心轨迹为（\quad）

\choices
    {抛物线}
    {双曲线}
    {椭圆}
    {圆}
\end{problem}

\begin{answer}
A
\end{answer}

\begin{solution}
设圆 \(C\) 的圆心为 \(P(x,y)\)，半径为 \(r\). 由于圆 \(C\) 与 \(y=0\) 相切，\(r=|y|\). 已知圆的圆心为 \((0,3)\)，半径为 \(1\)，外切条件给出
\[
x^2+(y-3)^2=(|y|+1)^2
\]
当 \(y\geqslant0\) 时，化简得 \(x^2-8y+8=0\)，即 \(y=\frac{x^2}{8}+1\). 这是一条抛物线，并且其所有点都满足 \(y\geqslant1\).

当 \(y<0\) 时，化简得 \(x^2-4y+8=0\)，即 \(y=\frac{x^2+8}{4}>0\)，与 \(y<0\) 矛盾. 因此圆心轨迹为抛物线，选 A.
\end{solution}

\begin{problem}
如图，某几何体的正视图（主视图）、侧视图（左视图）和俯视图分别是等边三角形、等腰三角形和菱形，则该几何体的体积为（\quad）

\bitmapfigure[max width=0.65\linewidth]{img/2011/guangdong_liberal/q9_fig1.png}

\choices
    {\(4\sqrt{3}\)}
    {\(4\)}
    {\(2 \sqrt{3}\)}
    {\(2\)}
\end{problem}

\begin{answer}
C
\end{answer}

\begin{solution}
由三视图可知，该几何体是以菱形为底面的四棱锥. 正视图为边长 \(2\sqrt{3}\) 的等边三角形，所以它的底边投影，即菱形的一条对角线，长为 \(2\sqrt{3}\)，四棱锥的高为
\[
h=\sqrt{(2\sqrt{3})^2-(\sqrt{3})^2}=3
\]
侧视图的底边给出菱形的另一条对角线，长为 \(2\). 这也与俯视图中边长为 \(2\) 的菱形相符，因为两条半对角线长为 \(\sqrt{3}\) 和 \(1\). 因此底面积为
\[
S=\frac{1}{2}\cdot2\sqrt{3}\cdot2=2\sqrt{3}
\]
所以体积为 \(V=\frac{1}{3}Sh=\frac{1}{3}\cdot2\sqrt{3}\cdot3=2\sqrt{3}\)，故选 C.
\end{solution}

\begin{problem}
设 \(f(x)\)，\(g(x)\)，\(h(x)\) 是 \(\R\) 上的任意实值函数. 如下定义两个函数 \((f \circ g)(x)\) 和 \((f \bullet g)(x)\)：对任意 \(x \in \R\)，\((f \circ g)(x) = f(g(x))\)；\((f \bullet g)(x) = f(x)g(x)\)，则下列等式恒成立的是（\quad）

\choices
    {\(\big((f\circ g)\bullet h\big)(x) = \big((f\bullet h)\circ (g\bullet h)\big)(x)\)}
    {\(\big((f\bullet g)\circ h\big)(x) = \big((f\circ h)\bullet (g\circ h)\big)(x)\)}
    {\(\big((f\circ g)\circ h\big)(x) = \big((f\circ h)\circ (g\circ h)\big)(x)\)}
    {\(\big((f\bullet g)\bullet h\big)(x) = \big((f\bullet h)\bullet (g\bullet h)\big)(x)\)}
\end{problem}

\begin{answer}
B
\end{answer}

\begin{solution}
对选项 B，左边为
\[
\big((f\bullet g)\circ h\big)(x)=(f\bullet g)(h(x))=f(h(x))g(h(x))
\]
右边为 \(\big((f\circ h)\bullet(g\circ h)\big)(x)=f(h(x))g(h(x))\)，所以 B 对任意 \(f\)，\(g\)，\(h\) 恒成立.

其余各项可分别构造反例. 对 A 取 \(f(x)=x\)，\(g(x)=1\)，\(h(x)=2\)，则左边为 \(2\)，右边为 \(4\)；对 C 取 \(f(x)=x\)，\(g(x)=x\)，\(h(x)=x+1\)，则左边为 \(x+1\)，右边为 \(x+2\)；对 D 取 \(f(x)=g(x)=h(x)=2\)，则左边为 \(8\)，右边为 \(16\). 因此只有 B 恒成立.
\end{solution}

\section{填空题}

\begin{problem}
已知 \(\{a_{n}\}\) 是递增的等比数列，\(a_{2} = 2\)，\(a_{4} - a_{3} = 4\). 则此数列的公比 \(q = \)\fillinblank{}.
\end{problem}

\begin{answer}
\(2\)
\end{answer}

\begin{solution}
设公比为 \(q\). 由 \(a_2=2\)，得 \(a_3=2q\)，\(a_4=2q^2\). 代入 \(a_4-a_3=4\)，得 \(2q^2-2q=4\)，即
\[
q^2-q-2=(q-2)(q+1)=0
\]
所以 \(q=2\) 或 \(q=-1\). 当 \(q=-1\) 时，数列项正负交替，不可能递增，故 \(q=2\).
\end{solution}

\begin{problem}
设函数 \( f(x) = x^{3}\cos x + 1 \). 若 \( f(a) = 11 \)，则 \( f(-a) = \)\fillinblank{}.
\end{problem}

\begin{answer}
\(-9\)
\end{answer}

\begin{solution}
由 \(f(a)=11\) 得 \(a^3\cos a+1=11\)，即 \(a^3\cos a=10\). 又因为 \(\cos(-a)=\cos a\)，所以
\[
f(-a)=(-a)^3\cos(-a)+1=-a^3\cos a+1=-10+1=-9
\]
\end{solution}

\begin{problem}
为了解篮球爱好者小李的投篮命中率与打篮球时间之间的关系. 下表记录了小李某月 \(1\) 号到 \(5\) 号每天打篮球时间 \(x\)（单位：小时）与当天投篮命中率 \(y\) 之间的关系：

\begin{center}
\begin{tabular}{c|ccccc}
时间 \(x\) & \(1\) & \(2\) & \(3\) & \(4\) & \(5\) \\
\hline
命中率 \(y\) & \(0.4\) & \(0.5\) & \(0.6\) & \(0.6\) & \(0.4\)
\end{tabular}
\end{center}

小李这 \(5\) 天的平均投篮命中率为\fillinblank{}；用线性回归分析的方法，预测小李该月 \(6\) 号打 \(6\) 小时篮球的投篮命中率为\fillinblank{}.
\end{problem}

\begin{answer}
\(0.5\)；\(0.53\)
\end{answer}

\begin{solution}
平均命中率为
\[
\bar y=\frac{0.4+0.5+0.6+0.6+0.4}{5}=0.5
\]
而 \(\bar x=\frac{1+2+3+4+5}{5}=3\). 线性回归直线 \(\hat y=\hat a+\hat b x\) 中，
\[
\hat b=\frac{\sum_{i=1}^{5}(x_i-\bar x)(y_i-\bar y)}{\sum_{i=1}^{5}(x_i-\bar x)^2}
 =\frac{(-2)(-0.1)+(-1)\cdot0+0\cdot0.1+1\cdot0.1+2(-0.1)}{(-2)^2+(-1)^2+0^2+1^2+2^2}
 =\frac{0.1}{10}=0.01
\]
因此 \(\hat a=\bar y-\hat b\bar x=0.5-0.01\times3=0.47\)，回归直线为 \(\hat y=0.47+0.01x\). 当 \(x=6\) 时，预测命中率为 \(\hat y=0.47+0.01\times6=0.53\).
\end{solution}

\section{选考题：填空题}

\begin{problem}
已知两曲线参数方程分别为 \(\begin{cases}
x = \sqrt{5} \cos \theta,\\
y = \sin \theta,
\end{cases}\) （\(0 \leqslant \theta < \pi\)）和 \(\begin{cases}
x = \frac{5}{4} t^2,\\
y = t,
\end{cases}\) （\(t \in \R\)），它们的交点坐标为\fillinblank{}.
\end{problem}

\begin{answer}
\(\left(1,\dfrac{2\sqrt{5}}{5}\right)\)
\end{answer}

\begin{solution}
第一条曲线满足 \(\frac{x^2}{5}+y^2=1\). 因为 \(0\leqslant\theta<\pi\)，有 \(y=\sin\theta\geqslant0\). 第二条曲线中 \(y=t\)，所以 \(x=\frac{5}{4}y^2\). 代入椭圆方程得
\[
\frac{1}{5}\left(\frac{5}{4}y^2\right)^2+y^2=1
\]
即 \(5y^4+16y^2-16=0\)，从而
\[
(5y^2-4)(y^2+4)=0
\]
由于 \(y\geqslant0\)，只能有 \(y^2=\frac45\)，故 \(y=\frac{2\sqrt5}{5}\)，再由 \(x=\frac54y^2\) 得 \(x=1\). 此时 \(x>0\)，\(y>0\)，可取第一象限的 \(\theta\in[0,\frac\pi2)\)，确实满足第一条参数方程；第二条取 \(t=\frac{2\sqrt5}{5}\) 即可. 因此交点坐标为 \(\left(1,\frac{2\sqrt5}{5}\right)\).
\end{solution}

\begin{problem}
如图，在梯形 \(ABCD\) 中，\(AB \parallel CD\)，\(AB = 4\)，\(CD = 2\)，\(E\)，\(F\) 分别为 \(AD\)，\(BC\) 上的点，且 \(EF = 3\)，\(EF \parallel AB\)，则梯形 \(ABFE\) 与梯形 \(EFCD\) 的面积比为\fillinblank{}.

\bitmapfigure[max width=0.36\linewidth]{img/2011/guangdong_liberal/q15_fig1.png}
\end{problem}

\begin{answer}
\(7:5\)
\end{answer}

\begin{solution}
设梯形 \(ABCD\) 的高为 \(h\)，并令 \(\frac{AE}{AD}=\frac{BF}{BC}=r\). 因为 \(AB\)，\(EF\)，\(CD\) 互相平行，平行线截两腰所得的对应线段成比例，且截线长度随位置线性变化，所以
\[
EF=(1-r)AB+rCD=4-2r
\]
由 \(EF=3\) 得 \(r=\frac12\)，故 \(EF\) 到两底的距离都为 \(\frac h2\). 于是
\(S_{ABFE}=\frac{4+3}{2}\cdot\frac h2=\frac{7h}{4}\)，\(S_{EFCD}=\frac{3+2}{2}\cdot\frac h2=\frac{5h}{4}\).
所求面积比为 \(7:5\).
\end{solution}

\section{解答题}

\begin{problem}
已知函数 \(f(x) = 2\sin \left(\frac{1}{3} x - \frac{\pi}{6}\right)\)，\(x \in \R\).
\begin{enumerate}
\item 求 \( f(0) \) 的值；
\item 设 \(\alpha\)，\(\beta \in \left[0, \frac{\pi}{2}\right]\)，\(f\left(3\alpha + \frac{\pi}{2}\right) = \frac{10}{13}\)，\(f(3\beta + 2\pi) = \frac{6}{5}\)，求 \(\sin (\alpha + \beta)\) 的值.
\end{enumerate}
\end{problem}

\begin{solution}
\begin{enumerate}
\item 代入 \(x=0\)，得 \(f(0)=2\sin\left(-\frac\pi6\right)=-1\).
\item 将 \(x=3\alpha+\frac\pi2\) 代入，得
\[
f\left(3\alpha+\frac\pi2\right)=2\sin\left(\alpha+\frac\pi6-\frac\pi6\right)=2\sin\alpha=\frac{10}{13}
\]
所以 \(\sin\alpha=\frac5{13}\). 由 \(\alpha\in[0,\frac\pi2]\)，得 \(\cos\alpha=\frac{12}{13}\). 同理，
\[
f(3\beta+2\pi)=2\sin\left(\beta+\frac{2\pi}{3}-\frac\pi6\right)=2\sin\left(\beta+\frac\pi2\right)=2\cos\beta=\frac65
\]
从而 \(\cos\beta=\frac35\). 由 \(\beta\in[0,\frac\pi2]\)，得 \(\sin\beta=\frac45\). 因此
\[
\sin(\alpha+\beta)=\sin\alpha\cos\beta+\cos\alpha\sin\beta
 =\frac5{13}\cdot\frac35+\frac{12}{13}\cdot\frac45=\frac{63}{65}
\]
\end{enumerate}
\end{solution}

\begin{problem}
在某次测验中，有 \(6\) 位同学的平均成绩为 \(75\) 分. 用 \( x_{n} \) 表示编号为 \(n\) （\(n = 1\)，\(2\)，\(\cdots\)，\(6\)）的同学所得成绩，且前 \(5\) 位同学的成绩如下：

\begin{center}
\begin{tabular}{c|ccccc}
编号 \(n\) & \(1\) & \(2\) & \(3\) & \(4\) & \(5\) \\
\hline
成绩 \(x_{n}\) & \(70\) & \(76\) & \(72\) & \(70\) & \(72\)
\end{tabular}
\end{center}

\begin{enumerate}
\item 求第 \(6\) 位同学的成绩 \(x_6\)，及这 \(6\) 位同学成绩的标准差 \(s\)；
\item 从前 \(5\) 位同学中，随机地选 \(2\) 位同学，求恰有 \(1\) 位同学成绩在区间 \((68, 75)\) 中的概率.
\end{enumerate}
\end{problem}

\begin{solution}
\begin{enumerate}
\item 六位同学的总成绩为 \(6\times75=450\)，前五位同学的总成绩为 \(70+76+72+70+72=360\)，所以 \(x_6=450-360=90\). 以平均成绩 \(\bar x=75\) 为基准，平方离差和为
\[
\sum_{i=1}^{6}(x_i-\bar x)^2=(-5)^2+1^2+(-3)^2+(-5)^2+(-3)^2+15^2=294
\]
因此标准差为 \(s=\sqrt{\frac{294}{6}}=7\).
\item 前五位同学中，成绩在 \((68,75)\) 内的有编号 \(1\)，\(3\)，\(4\)，\(5\) 的四位同学，成绩为 \(76\) 的二号同学不在该区间内. 随机选两位的基本事件数为 \(\mathrm C_5^2\)，恰有一位来自这四位同学的事件数为 \(\mathrm C_4^1\mathrm C_1^1\). 故所求概率为
\[
\frac{\mathrm C_4^1\mathrm C_1^1}{\mathrm C_5^2}=\frac4{10}=\frac25
\]
\end{enumerate}
\end{solution}

\begin{problem}
如图所示的几何体是将高为 \(2\)、底面半径为 \(1\) 的直圆柱沿过轴的平面切开后，将其中一半沿切面向右水平平移后得到的. \(A\)，\(A'\)，\(B\)，\(B'\) 分别为 \(\myarc{CD}\)，\(\myarc{C'D'}\)，\(\myarc{DE}\)，\(\myarc{D'E'}\) 的中点，\(O_1\)，\(O_1'\)，\(O_2\)，\(O_2'\) 分别为 \(CD\)，\(C'D'\)，\(DE\)，\(D'E'\) 的中点.

\begin{enumerate}
\item 证明：\(O_{1}^{\prime}\)，\(A^{\prime}\)，\(O_{2}\)，\(B\) 四点共面；
\item 设 \(G\) 为 \(A^{\prime}A\) 的中点，延长 \(A^{\prime}O_{1}^{\prime}\) 到 \(H^{\prime}\)，使得 \(O_{1}^{\prime}H^{\prime}=A^{\prime}O_{1}^{\prime}\)，证明：\(BO_{2}^{\prime}\perp\) 平面 \(H^{\prime}B^{\prime}G\).
\end{enumerate}

\bitmapfigure[max width=0.44\linewidth]{img/2011/guangdong_liberal/q18_fig1.png}
\end{problem}

\begin{solution}
建立直角坐标系：取 \(D\) 为原点，\(DE\) 所在方向为 \(x\) 轴，竖直向上为 \(z\) 轴，左侧半圆柱位于 \(y\geqslant0\) 一侧，右侧半圆柱位于 \(y\leqslant0\) 一侧. 由半径为 \(1\)、平移后两半圆柱的直径分别为 \(CD\) 和 \(DE\)，可取
\(C=(-2,0,0)\)，\(D=(0,0,0)\)，\(E=(2,0,0)\).

\(C'=(-2,0,2)\)，\(D'=(0,0,2)\)，\(E'=(2,0,2)\).
于是 \(O_1=(-1,0,0)\)，\(O_1'=(-1,0,2)\)，\(O_2=(1,0,0)\)，\(O_2'=(1,0,2)\).

以及 \(A=(-1,1,0)\)，\(A'=(-1,1,2)\)，\(B=(1,-1,0)\)，\(B'=(1,-1,2)\).

\begin{enumerate}
\item 由 \(O_1'=(-1,0,2)\)，\(A'=(-1,1,2)\)，\(O_2=(1,0,0)\) 三点可知平面 \(O_1'A'O_2\) 的方程为 \(x+z=1\). 而 \(B=(1,-1,0)\) 也满足 \(x+z=1\)，所以 \(O_1'\)，\(A'\)，\(O_2\)，\(B\) 四点共面.
\item 因为 \(G\) 是 \(A'A\) 的中点，所以 \(G=(-1,1,1)\). 从 \(A'\) 沿着 \(A'O_1'\) 的方向延长一段 \(A'O_1'\)，得
\[
H'=2O_1'-A'=(-1,-1,2)
\]
因此 \(\overrightarrow{BO_2'}=(0,1,2)\)，\(\overrightarrow{H'B'}=(2,0,0)\)，\(\overrightarrow{H'G}=(0,2,-1)\).
于是 \(\overrightarrow{BO_2'}\cdot\overrightarrow{H'B'}=0\)，且 \(\overrightarrow{BO_2'}\cdot\overrightarrow{H'G}=0\). 向量 \(\overrightarrow{H'B'}\) 与 \(\overrightarrow{H'G}\) 不共线，它们确定平面 \(H'B'G\)，故 \(BO_2'\perp\) 平面 \(H'B'G\).
\end{enumerate}
\end{solution}

\begin{problem}
设 \(a > 0\)，讨论函数 \(f(x) = \ln x + a(1 - a)x^2 - 2(1 - a)x\) 的单调性.
\end{problem}

\begin{solution}
函数定义域为 \((0,+\infty)\)，且
\[
f'(x)=\frac1x+2a(1-a)x-2(1-a)=\frac{g(x)}{x}
\]
其中 \(g(x)=2a(1-a)x^2-2(1-a)x+1\).
由于 \(x>0\)，\(f'(x)\) 与 \(g(x)\) 同号.

当 \(0<a<\frac13\) 时，\(g(x)\) 的判别式为 \(4(1-a)(1-3a)>0\)，两个根为
\(x_1=\frac{1-\sqrt{\frac{1-3a}{1-a}}}{2a}\)，\(x_2=\frac{1+\sqrt{\frac{1-3a}{1-a}}}{2a}\).
因为 \(0<\sqrt{\frac{1-3a}{1-a}}<1\)，所以 \(0<x_1<x_2\). 此时二次项系数为正，故 \(g(x)>0\) 在 \((0,x_1)\) 和 \((x_2,+\infty)\) 上成立，\(g(x)<0\) 在 \((x_1,x_2)\) 上成立. 因此 \(f(x)\) 在 \((0,x_1)\) 和 \((x_2,+\infty)\) 上单调递增，在 \((x_1,x_2)\) 上单调递减.

当 \(a=\frac13\) 时，
\[
g(x)=\frac49x^2-\frac43x+1=\left(\frac{2x}{3}-1\right)^2\geqslant0
\]
虽然仅在 \(x=\frac32\) 时导数为零，但函数在整个 \((0,+\infty)\) 上单调递增.

当 \(\frac13<a<1\) 时，二次项系数为正且判别式 \(4(1-a)(1-3a)<0\)，所以 \(g(x)>0\). 当 \(a=1\) 时，\(g(x)=1>0\). 因此 \(\frac13\leqslant a\leqslant1\) 时，\(f(x)\) 在 \((0,+\infty)\) 上单调递增.

当 \(a>1\) 时，令 \(c=a-1>0\)，则
\[
g(x)=-2acx^2+2cx+1
\]
其两个根分别为 \(\frac{1-\sqrt{\frac{3a-1}{a-1}}}{2a}<0\) 和 \(x_0=\frac{1+\sqrt{\frac{3a-1}{a-1}}}{2a}>0\).
因二次项系数为负，故在正半轴上 \(g(x)>0\)（\(0<x<x_0\)），\(g(x)<0\)（\(x>x_0\)）. 所以 \(f(x)\) 在 \((0,x_0)\) 上单调递增，在 \((x_0,+\infty)\) 上单调递减.
\end{solution}

\begin{problem}
设 \(b > 0\)，数列 \(\{a_{n}\}\) 满足 \(a_{1} = b\)，\(a_{n} = \frac{nba_{n - 1}}{a_{n - 1} + n - 1} \)（\(n \geqslant 2\)）.
\begin{enumerate}
\item 求数列 \( \{a_{n}\} \) 的通项公式；
\item 证明：对于一切正整数 \(n\)，\(2a_{n} \leqslant b^{n+1} + 1\).
\end{enumerate}
\end{problem}

\begin{solution}
\begin{enumerate}
\item 先证数列各项均为正. 已知 \(a_1=b>0\)，若 \(a_{n-1}>0\)，则递推式的分子、分母均为正，从而 \(a_n>0\).

令 \(S_n=1+b+b^2+\cdots+b^{n-1}\). 当 \(n=1\) 时，\(a_1=b=\frac{1\cdot b^1}{S_1}\). 假设 \(a_{n-1}=\frac{(n-1)b^{n-1}}{S_{n-1}}\)，则
\[
a_{n-1}+n-1=\frac{(n-1)b^{n-1}}{S_{n-1}}+(n-1)=\frac{(n-1)S_n}{S_{n-1}}
\]
代入递推式得
\[
a_n=\frac{nb\cdot\frac{(n-1)b^{n-1}}{S_{n-1}}}{\frac{(n-1)S_n}{S_{n-1}}}=\frac{nb^n}{S_n}
\]
所以对一切正整数 \(n\)，
\[
a_n=\frac{nb^n}{1+b+b^2+\cdots+b^{n-1}}
\]
\item 由均值不等式，
\[
b^{n+1}+1\geqslant2b^{\frac{n+1}{2}}
\]
又因 \(b>0\)，对 \(1\)，\(b\)，\(b^2\)，\(\ldots\)，\(b^{n-1}\) 使用均值不等式，得
\[
S_n\geqslant n\sqrt[n]{b^{0+1+\cdots+(n-1)}}=nb^{\frac{n-1}{2}}
\]
两式相乘，得到 \((b^{n+1}+1)S_n\geqslant2nb^n\). 由于 \(S_n>0\)，再利用 \(a_n=\frac{nb^n}{S_n}\)，可得 \(b^{n+1}+1\geqslant2a_n\)，即所证不等式成立.
\end{enumerate}
\end{solution}

\begin{problem}
在平面直角坐标系 \(xOy\) 中，直线 \(l: x = -2\) 交 \(x\) 轴于点 \(A\). 设 \(P\) 是 \(l\) 上一点，\(M\) 是线段 \(OP\) 的垂直平分线上一点，且满足 \(\angle MPO = \angle AOP\).
\begin{enumerate}
\item 当点 \(P\) 在 \(l\) 上运动时，求点 \(M\) 的轨迹 \(E\) 的方程；
\item 已知 \(T(1,-1)\)，设 \(H\) 是 \(E\) 上动点，求 \( |HO| + |HT| \) 的最小值，并给出此时点 \(H\) 的坐标；
\item 过点 \(T(1, -1)\) 且不平行于 \(y\) 轴的直线 \(l_{1}\) 与轨迹 E 有且只有两个不同的交点，求直线 \(l_{1}\) 的斜率 \(k\) 的取值范围.
\end{enumerate}
\end{problem}

\begin{solution}
\begin{enumerate}
\item 设 \(P=(-2,t)\)，\(M=(x,y)\). 因为 \(M\) 在线段 \(OP\) 的垂直平分线上，所以 \(MO=MP\)，从而
\[
x^2+y^2=(x+2)^2+(y-t)^2
\]
即
\[
4x+4-2ty+t^2=0
\]
在等腰三角形 \(MOP\) 中，\(MO=MP\)，所以 \(\angle MOP=\angle MPO=\angle AOP\). 因此，从射线 \(OP\) 看，射线 \(OM\) 有如下两种可能位置.

若射线 \(OM\) 与射线 \(OA\) 重合，则 \(y=0\). 代入上式得 \(x=-\frac{t^2+4}{4}\leqslant-1\)，且当 \(t\) 取遍实数时得到射线 \(y=0\)，\(x\leqslant-1\).

另一种情形是射线 \(OP\) 平分 \(\angle AOM\). 令 \(u=\overrightarrow{OP}=(-2,t)\)，\(v=\overrightarrow{OA}=(-2,0)\). 将 \(v\) 关于直线 \(OP\) 对称，所得方向向量为
\[
2\frac{v\cdot u}{u\cdot u}u-v
=2\frac{4}{t^2+4}(-2,t)-(-2,0)
=\frac{2}{t^2+4}(t^2-4,4t)
\]
故可设 \(M=\lambda(t^2-4,4t)\)，其中 \(\lambda>0\). 代入（1）得
\[
4\lambda(t^2-4)+4-8\lambda t^2+t^2=0
\]
所以 \(\lambda=\frac14\)，即
\[
M=\left(\frac{t^2-4}{4},t\right)
\]
消去参数得 \(y^2=4x+4\)，且此支满足 \(x\geqslant-1\). 综上，轨迹为 \(E\) 的两部分：\(y^2=4x+4\)（\(x\geqslant-1\)）以及 \(y=0\)（\(x\leqslant-1\)）.
两部分在点 \((-1,0)\) 相接.
\item 对抛物线支 \(y^2=4(x+1)\) 上的点 \(H=(x,y)\)，焦点为 \(O=(0,0)\)，准线为 \(d:x=-2\). 设 \(H_0=(-2,y)\) 为 \(H\) 在准线上的垂足. 由抛物线定义，\(HO=HH_0\)，故
\[
HO+HT=HH_0+HT\geqslant H_0T=\sqrt{9+(y+1)^2}\geqslant3
\]
当且仅当 \(y=-1\) 且 \(H_0\)，\(H\)，\(T\) 共线时取等号. 此时由抛物线方程得 \(H=\left(-\frac34,-1\right)\)，并且该点确实在抛物线上.

对另一支 \(y=0\)，\(x\leqslant-1\)，令 \(u=-x\geqslant1\)，则
\[
HO+HT=u+\sqrt{(u+1)^2+1}>u+(u+1)\geqslant3
\]
在端点 \(u=1\) 时不等式仍为严格不等式. 因此全轨迹上的最小值为 \(3\)，此时 \(H=\left(-\frac34,-1\right)\).
\item 设直线 \(l_1\) 的斜率为 \(k\). 因其不平行于 \(y\) 轴，方程可写为 \(y+1=k(x-1)\).

先考虑 \(k\ne0\). 由 \(x=1+\frac{y+1}{k}\)，代入抛物线方程得
\[
ky^2-4y-8k-4=0
\]
其判别式为
\[
\Delta=(-4)^2-4k(-8k-4)=16(2k^2+k+1)
\]
由于 \(2k^2+k+1>0\)，直线与抛物线恒有两个不同的交点.

直线与附加支 \(y=0\)，\(x\leqslant-1\) 的交点横坐标为 \(x=1+\frac1k\). 当且仅当 \(-\frac12<k<0\) 时，\(1+\frac1k<-1\)，从而会在抛物线的两个交点之外再增加一个交点. 当 \(k=-\frac12\) 时，附加支交点为 \((-1,0)\)，它正好是抛物线顶点，不增加新的交点.

当 \(k=0\) 时，直线为 \(y=-1\)，与抛物线只有一个交点，故不符合题意. 因此恰有两个不同交点时
\[
k\in\left(-\infty,-\frac12\right]\cup(0,+\infty)
\]
\end{enumerate}
\end{solution}
